% TOP_PROBLEM: {"id": 1101209, "title": "Questions in Geometric Group Theory — Q 12.9", "category_id": 17, "category": {"id": 17, "name": "group_theory", "display_name": "Group Theory", "description": "Problems about groups, group actions, representations, and related algebraic structures.", "slug": "group-theory", "order_index": 17, "created_at": "2026-07-31T15:26:25.671Z"}, "status": "partially_solved", "problem_number": "AMR-010-1209", "set_id": 13}
% FIRST_POSTED: 2026-10-01
% ENTRY_KIND: statement_only
% UnsolvedMath ID 1101209; source catalogue code AMR-010-1209
% !TeX program = lualatex
% Definition and sources only; this file is not an LLM solution attempt.
\documentclass[11pt,a4paper]{article}
\usepackage[margin=25mm]{geometry}
\usepackage{fontspec}
\setmainfont{lmroman10-regular.otf}[BoldFont=lmroman10-bold.otf,
 ItalicFont=lmroman10-italic.otf,BoldItalicFont=lmroman10-bolditalic.otf]
\usepackage{amsmath,amssymb,mathtools}
\usepackage{unicode-math}
\setmathfont{latinmodern-math.otf}
\AtBeginDocument{\let\setminus\smallsetminus}
\usepackage{xurl,hyperref}
\hypersetup{colorlinks=true,urlcolor=blue}
\setcounter{secnumdepth}{0}
\setlength{\parindent}{0pt}
\setlength{\parskip}{5pt}
\setlength{\emergencystretch}{3em}
\begin{document}
\section*{Questions in Geometric Group Theory — Q 12.9}\label{1101209}
{\small UnsolvedMath ID 1101209 · AMR-010-1209 · Group Theory · AMR Open Problem Lists}
\subsection{Definitions and mathematical statement}
For a closed orientable surface $\Sigma_g$, $g\geq2$, let $\operatorname{Mod}(\Sigma_g)$ be the group of orientation-preserving homeomorphisms modulo isotopy. For a surface $\Sigma_{h,p}$ of genus $h\geq0$ with $p\geq1$ punctures and no boundary, use the analogous mapping class group in which punctures may be permuted. Isotopies preserve the puncture set. The subgroup fixing punctures individually has finite index, so the virtual versions are unaffected by this convention.

Let $F_n$ be a finite-rank free group and put $\operatorname{Out}(F_n)=\operatorname{Aut}(F_n)/\operatorname{Inn}(F_n)$. For a given genus $g$, ask whether there is an injective homomorphism
\[
 \operatorname{Mod}(\Sigma_g)\longrightarrow\operatorname{Out}(F_n)
\]
for some finite $n$, or at least such an injection from a finite-index subgroup of $\operatorname{Mod}(\Sigma_g)$. Ask the same two questions with target $\operatorname{Mod}(\Sigma_{h,p})$, allowing $h$ and positive $p$ to vary.

The problem seeks to determine which genera and targets permit these embeddings. Both the full-group and virtual variants are retained. No compatibility with homology representations, a covering map between surfaces, or a specified free-group action is assumed. An injection into a mapping class group of another closed surface does not answer the punctured-target question. Likewise, a homomorphism having finite or nontrivial image but nontrivial kernel is not an embedding. The punctures are essential distinguished missing points rather than extra unmarked points.
\subsection{Short English statement}
Can closed-surface mapping class groups, or finite-index subgroups of them, embed into $\operatorname{Out}(F_n)$ or mapping class groups of punctured surfaces?
\subsection{Sources}
\begin{itemize}
\item[S1] ULAM.AI and the UnsolvedMath Contributors, supplied problems.json, record 1101209 (AMR-010-1209), AMR Open Problem Lists. Catalogue identity and source transcription; CC BY 4.0.
\url{https://huggingface.co/datasets/ulamai/UnsolvedMath}.
\item[S2] Bestvina - Questions in Geometric Group Theory (pdf) (2004), Question 12.9 (PDF page 21). Original source indexed by AMR-010-1209.
\url{https://www.math.utah.edu/~bestvina/eprints/questions-updated.pdf}.
\end{itemize}
\end{document}
